\section{Classification comes before the loan calculation}
\label{sec:bank-asset-classification}

A loan's contractual principal is not enough to determine its carrying
amount. The dated amortised-cost examples that follow require both an
eligible contractual cash-flow pattern and a business model for collecting
those flows. These are findings about the instrument and the way assets are
managed, not labels chosen to obtain a preferred profit figure. The following
procedure uses the retained historical IFRS~9 clauses in EU
Regulation~2023/1803; the surrounding Australian and Chinese examples retain
their own stated jurisdiction and date boundaries.

First establish the financial asset's terms and the level at which the bank
actually manages it. Evidence includes how performance is reported to key
management, how risks are managed and how managers are compensated. An
intention to retain one asset is insufficient. Sales must be explained by
their reasons, timing, volume and place in that management objective; the
source supplies no universal sale-count threshold
\citep[IFRS~9.4.1.1--3, B4.1.1--4B]{eu20231803}.

Next ask whether contractual payments are solely principal and interest on
the outstanding principal. Principal here begins with fair value at initial
recognition. Interest compensates for time value, credit risk, other basic
lending risks and costs, and a profit margin. Equity-linked returns and
leverage introduce different exposures. For a hybrid whose host is a
financial asset, assess the entire contract under IFRS~9.4.3.2; the issuer's
separate liability and equity components do not automatically become the
holder's measurement units \citep[IFRS~9.4.1.3, 4.3.2, B4.1.7A--9,
B4.1.14]{eu20231803}.

If that cash-flow test passes, a hold-to-collect business model leads to
amortised cost, while a model achieved by both collection and sales leads to
fair value through other comprehensive income (FVOCI) for debt. Other
financial assets generally use fair value through profit or loss (FVTPL),
subject to the separate equity election below. At initial recognition,
an otherwise eligible asset may be irrevocably designated FVTPL when the
specified accounting-mismatch condition is met. Preserve that designation
and its reason. A later change in one asset's intended sale date is not by
itself a business-model change authorising reclassification
\citep[IFRS~9.4.1.2--5, 4.4.1, B4.1.2A--3]{eu20231803}.

\subsection{The same price decline can reach profit and OCI differently}

Consider a debt asset purchased for 100 and now worth 98. Supply cumulative
ECL of 3, with no interest accrual, foreign-exchange movement, tax or fees
in this measurement comparison. Opening cash and equity are each 200;
purchase leaves cash 100. The three cases below have different admitted
business models or an eligible FVTPL designation. They compare reporting
consequences, rather than three freely selectable policies for one factual
case.

At amortised cost, record impairment expense and allowance 3. Gross asset
100 less allowance 3 gives net asset 97, so cash 100 plus net asset 97 equals
equity 197. A subsequent derecognising sale for 98 debits cash 98 and
allowance 3, credits gross asset 100 and recognises gain 1. Cumulative
profit is \(-3+1=-2\), and closing cash and equity are 198.

For debt FVOCI, the balance-sheet asset is 98. Record the fair-value decline
2 in OCI. Impairment expense 3 enters profit, with its allowance counterpart
credited to OCI rather than deducted again from the asset's fair value.
Net OCI is therefore \(-2+3=1\); profit is \(-3\), and total
comprehensive income is \(-2\). Cash 100 plus asset 98 equals equity
198. Reporting asset 95 by subtracting the allowance once more would count
the credit effect twice. On a derecognising sale for 98, remove the asset
and reclassify cumulative OCI gain 1 to profit. Cumulative profit then
equals \(-2\), with OCI cleared
\citep[IFRS~9.5.5.1--2, 5.7.10--11]{eu20231803}.

Under FVTPL, the decline from 100 to 98 is expense 2. There is no separate
IFRS~9 loss allowance added to that fair-value loss. The sale at the same
price creates no further gain. All three examples ultimately close with
cash and equity 198, but their interim carrying amounts and placement of
income differ. Reconciliation must preserve the chosen route rather than
infer it from the equal final cash.

An equity investment is a distinct instrument. Where the irrevocable FVOCI
election is available and made at initial recognition, a fall from 100 to
98 gives OCI loss 2 and no fair-value expense in profit. It excludes an
investment held for trading and specified contingent consideration of an
acquirer in a business combination. On sale, the cumulative loss is not
recycled to profit; a transfer within equity is permitted. Dividends follow
their separate recognition conditions, including whether they clearly
represent recovery of part of the investment's cost. With no dividends in
this example, cash and equity again end at 198 while cumulative profit is
zero \citep[IFRS~9.5.7.5--6, B5.7.1--1A]{eu20231803}. Copying the debt
recycling entry into this case would wrongly create expense 2 on sale.

\subsection{A familiar interest label can conceal a different exposure}

A floating rate need not faithfully represent the passage of time. Suppose
principal 100 pays monthly interest but resets each month to a one-year
rate. A comparison instrument would reset monthly to a one-month rate,
with the other relevant terms and credit quality matched. In one supplied
six-month scenario, hold the contractual annualised rate at 12\% and the
benchmark rate at 6\%. Each month's interest is respectively 1 and 0.5;
six payments differ by 3. With principal repaid at month six, total
undiscounted receipts are 106 and 103.

This calculation identifies the difference to investigate. IFRS~9's
modified-time-value test compares undiscounted contractual and benchmark
cash flows both by reporting period and cumulatively over the instrument's
life, considering reasonably possible scenarios. A current yield curve,
discounted present-value match or an invented percentage cutoff cannot
settle whether the difference could be significant. Retain the benchmark
construction, scenarios, period totals and significance reasoning before
making that finding \citep[IFRS~9.B4.1.9A--9E]{eu20231803}.

Nonrecourse terms likewise do not automatically fail the cash-flow test.
The holder examines the underlying assets or cash flows to determine whether
the contractual restriction introduces payments inconsistent with basic
lending. Contractually linked tranches require distinct tests of the
tranche's terms, underlying pool and relative credit-risk concentration,
including the permitted instruments in the pool. If the holder cannot
perform the assessment required by B4.1.26, that clause requires FVTPL.
A missing document in this system, however, is initially an unresolved
evidence requirement; it is not proof that the holder made, or could not
make, that assessment \citep[IFRS~9.B4.1.15--26]{eu20231803}.

The method therefore records the governing rule, management evidence,
contractual findings, designation or election, and remaining unknowns before
selecting an interest, impairment and disposal procedure. Unsupported
classification leaves those dependent calculations conditional. A supplied
fair value or loss allowance can illustrate the reporting arithmetic while
its valuation or credit model still requires separate specification and
validation.
